\scanpage{166}
\mspage{115}
% archive sheet number: 166
2.5. Complément : \deleted{\unsure{intégration} d'un préfaisceau} \inserted{critère d'asphéricité totale} \deleted{sur}\inserted{d'}une catégorie \inserted{$W$-}propre sur une autre.

Soit

\[
f:X\longrightarrow I
\]
une catégorie propre sur une catégorie d'indices $I$. Si $I$ a un objet final $i_0$, alors \deleted{\unsure{le}} $X_{/i_0}$ s'identifie à $X$ et par suite l'inclusion \deleted{$X_{i_0}\longrightarrow X$ de}

\[
X_{i_0}\unsure{-|\longrightarrow}X
\]
de la fibre de $X$ en $i_0$ est $W$-asphérique, et \unsure{a|à} fortiori est dans $W$. Donc $X$ est asphérique \unsure{ssi|sss} $X_{i_0}$ l'est. Nous voulons \unsure{étendre|élever} ce résultat en un critère de \uline{totale} $W$-asphéricité, de façon précise prouver la

\uline{Proposition} 11 : Soit \deleted{\unsure{$I$|$f$}}\unsure{$f$}$:X-I$ $W$-propre, et \margin{$\searrow\!\searrow$}soit $i_0$ un objet final de $I$. Si $X_{i_0}$ est tot.\textsuperscript{t} $W$-asphérique, \inserted{(déf.~1, page \unsure{88|89})} il en est de même de $X$. (Sous réserve que $W$ satisfasse l'axiome de connexité $W(5)$.)

La démonstration sera donnée plus bas, après quelques préliminaires ayant un intérêt indépendant. Tout d'abord, rappelons que si $X$ est \inserted{$W$-}propre sur $I$, alors il donne naissance à un foncteur

\[
\begin{aligned}
I&\longrightarrow\mathrm{Hot}_W\\
i&\longmapsto\mathrm{hot}_W(X_i)\quad(\simeq\mathrm{hot}_W(X_{/i}))\,;
\end{aligned}
\]
en le composant avec le foncteur

\[
\pi_0:\mathrm{Hot}_W-\mathrm{Ens}
\]
(sous réserve que l'axiome \inserted{de connexité} $W(5)$, \deleted{\unsure{ait}} (p.~10) soit satisfait,
